% =====================================================================
% PROPOSED METHODOLOGY  ---  RESPONSABLE: INTEGRANTE 2
% =====================================================================
\section{Proposed Methodology}
\label{sec:methodology}

\subsection{Problem Formalization}
% TODO Int.2: notacion. Una carrera r es un grupo de n_r caballos.
% Definir si se modela P(gana | x) por caballo y luego se normaliza
% dentro de la carrera, o si se usa un modelo de ranking directo.
% \begin{equation} ... \end{equation}

\subsection{Proposed Pipeline}
% TODO Int.2: datos crudos -> limpieza -> feature engineering historico
% -> split temporal -> entrenamiento -> calibracion -> ranking ->
% evaluacion.

\subsection{Experimental Protocol}
% TODO Int.2 (coordinar con Int.3):
%  - Particion TEMPORAL y agrupada por race_id (no aleatoria).
%  - Validacion respetando el orden del tiempo.
%  - Metricas Y SU JUSTIFICACION: top-1 acc., MRR/NDCG a nivel de
%    carrera; log loss y Brier para calibracion.

\subsection{Baseline Plan}
% TODO Int.2: minimo tres.
%  (a) heuristico (favorito por peso o posicion de salida)
%  (b) regresion logistica simple
%  (c) el numero de la literatura (Seccion IV, Int.3)

\subsection{Risks and Mitigation}
% TODO Int.2: tabla corta. Leakage, desbalance, cambio de distribucion
% 1990-2020, campos de tamano variable, valores faltantes.
